% GENERATED FILE. Edit canonical lessons, then run npm run generate:books.
% canonical-source-hash: fnv1a64:d8156f53d28b1196
% canonical-lessons: JA-C351-man-ichi, JA-C351-hoken, JA-C351-kakunin-suru, JA-C351-anpi, JA-C351-nen-no-tame, JA-R351-review-347-351

\chapter{In Case Something Happens}
\label{ch:ja-in-case-something-happens}
\hlchaptermodality{351}

\begin{chapteropening}
\textbf{By the end of this chapter:} \emph{I can say what to keep ready in case something goes wrong.}

Complete the last lesson of chapter 351: I can say what to keep ready in case something goes wrong.
\end{chapteropening}

\section[\texorpdfstring{\ja{まんいち} (man'ichi)}{まんいち (man'ichi)}]{\ja{まんいち} (man'ichi) — in the unlikely event; by any chance}

\label{lesson:JA-C351-man-ichi}



\begin{quote}
Before the new one: say the Japanese for a countermeasure, then the Japanese for disaster preparedness.
\end{quote}

\subsection*{You'll want to know: \ja{まんいち}}
\textbf{\ja{まんいち}} — \emph{man'ichi} — \textquotedblleft{}in the unlikely event; by any chance\textquotedblright{}.

\begin{quote}
\textbf{\ja{まんいち} \ja{おくれたら、でんわして} \ja{ください。}} — \emph{man'ichi okuretara, denwa shite kudasai} — \textquotedblleft{}If by any chance you are late, please phone.\textquotedblright{}
\end{quote}

\begin{cousinweb}
\textbf{\ja{まんいち}} (\emph{man'ichi}) is Sino-Japanese (kango): two Chinese characters read with their on'yomi, readings borrowed from Chinese. They are \emph{man}, \textquotedblleft{}ten thousand\textquotedblright{} (Middle Chinese \emph{mjonH}) and \emph{ichi}, \textquotedblleft{}one\textquotedblright{} (\emph{'jit}) (KANJIDIC2; Middle Chinese in Baxter's transcription).
\end{cousinweb}

\begin{grammarlens}[title={What you've built}]
One more word for saying what happens if.
\end{grammarlens}

\subsection*{Your turn}
\begin{itemize}
\raggedright
  \item \emph{Say it:} \emph{man'ichi}
  \item \emph{Say it:} \emph{man'ichi}, once more
  \item \emph{Say it:} \emph{man'ichi okuretara, denwa shite kudasai} — If by any chance you are late, please phone.
  \item \emph{Recall:} say the Japanese for a countermeasure, then the Japanese for disaster preparedness, then say \emph{man'ichi} again
\end{itemize}

\begin{tcolorbox}[breakable,colback=teal!4,colframe=teal!35!black,title={Before you move on}]
What does \ja{まんいち} mean? (\textquotedblleft{}In the unlikely event; by any chance\textquotedblright{}.) Say it once more.
\end{tcolorbox}
\section[\texorpdfstring{\ja{ほけん} (hoken)}{ほけん (hoken)}]{\ja{ほけん} (hoken) — insurance}

\label{lesson:JA-C351-hoken}



\begin{quote}
Before the new one: say the Japanese for disaster preparedness, then the Japanese for in the unlikely event.
\end{quote}

\subsection*{You'll want to know: \ja{ほけん}}
\textbf{\ja{ほけん}} — \emph{hoken} — \textquotedblleft{}insurance\textquotedblright{}.

\begin{quote}
\textbf{\ja{ほけんが} \ja{あったら、あんしんでしょう。}} — \emph{hoken ga attara, anshin deshō} — \textquotedblleft{}If you have insurance, you can feel easy.\textquotedblright{}
\end{quote}

\begin{cousinweb}
\textbf{\ja{ほけん}} (\emph{hoken}) is Sino-Japanese (kango): two Chinese characters read with their on'yomi, readings borrowed from Chinese. They are \emph{ho}, \textquotedblleft{}preserve\textquotedblright{} (Middle Chinese \emph{pawX}) and \emph{ken}, \textquotedblleft{}steep place\textquotedblright{} (\emph{xjemX}) (KANJIDIC2; Middle Chinese in Baxter's transcription).
\end{cousinweb}

\begin{grammarlens}[title={What you've built}]
One more word for saying what happens if.
\end{grammarlens}

\subsection*{Your turn}
\begin{itemize}
\raggedright
  \item \emph{Say it:} \emph{hoken}
  \item \emph{Say it:} \emph{hoken}, once more
  \item \emph{Say it:} \emph{hoken ga attara, anshin deshō} — If you have insurance, you can feel easy.
  \item \emph{Recall:} say the Japanese for disaster preparedness, then the Japanese for in the unlikely event, then say \emph{hoken} again
\end{itemize}

\begin{tcolorbox}[breakable,colback=teal!4,colframe=teal!35!black,title={Before you move on}]
What does \ja{ほけん} mean? (\textquotedblleft{}Insurance\textquotedblright{}.) Say it once more.
\end{tcolorbox}
\section[\texorpdfstring{\ja{かくにんする} (kakunin suru)}{かくにんする (kakunin suru)}]{\ja{かくにんする} (kakunin suru) — to confirm; to check}

\label{lesson:JA-C351-kakunin-suru}



\begin{quote}
Before the new one: say the Japanese for in the unlikely event, then the Japanese for insurance.
\end{quote}

\subsection*{You'll want to know: \ja{かくにんする}}
\textbf{\ja{かくにんする}} — \emph{kakunin suru} — \textquotedblleft{}to confirm; to check\textquotedblright{}.

\begin{quote}
\textbf{\ja{まちがえたら、もう} \ja{いちど} \ja{かくにんして} \ja{ください。}} — \emph{machigaetara, mō ichido kakunin shite kudasai} — \textquotedblleft{}If you make a mistake, please check it again.\textquotedblright{}
\end{quote}

Polite \textbf{\ja{かくにんします}}, \emph{kakuninshimasu}; a \ja{する} verb. Past (polite): \textbf{\ja{てんいんが} \ja{よやくを} \ja{かくにんしました。}}, \emph{ten'in ga yoyaku o kakunin shimashita}.

\begin{cousinweb}
\textbf{\ja{かくにんする}} (\emph{kakunin suru}) is built from a Sino-Japanese noun and the native verb \emph{suru}, \textquotedblleft{}to do\textquotedblright{}. The noun, \emph{kakunin}, is two Chinese characters read with their on'yomi, readings borrowed from Chinese: \emph{kaku}, \textquotedblleft{}firm\textquotedblright{} (Middle Chinese \emph{khaewk}) and \emph{nin}, \textquotedblleft{}recognize\textquotedblright{} (\emph{nyinH}) (KANJIDIC2; JMdict; Middle Chinese in Baxter's transcription).
\end{cousinweb}

\begin{grammarlens}[title={What you've built}]
One more word for saying what happens if.
\end{grammarlens}

\subsection*{Your turn}
\begin{itemize}
\raggedright
  \item \emph{Say it:} \emph{kakunin suru}
  \item \emph{Say it:} \emph{kakunin suru}, once more
  \item \emph{Say it:} \emph{machigaetara, mō ichido kakunin shite kudasai} — If you make a mistake, please check it again.
  \item \emph{Recall:} say the Japanese for in the unlikely event, then the Japanese for insurance, then say \emph{kakunin suru} again
\end{itemize}

\begin{tcolorbox}[breakable,colback=teal!4,colframe=teal!35!black,title={Before you move on}]
What does \ja{かくにんする} mean? (\textquotedblleft{}To confirm; to check\textquotedblright{}.) Say it once more.
\end{tcolorbox}
\section[\texorpdfstring{\ja{あんぴ} (anpi)}{あんぴ (anpi)}]{\ja{あんぴ} (anpi) — whether someone is safe}

\label{lesson:JA-C351-anpi}



\begin{quote}
Before the new one: say the Japanese for insurance, then the Japanese for to confirm.
\end{quote}

\subsection*{You'll want to know: \ja{あんぴ}}
\textbf{\ja{あんぴ}} — \emph{anpi} — \textquotedblleft{}whether someone is safe\textquotedblright{}.

\begin{quote}
\textbf{\ja{たいふうが} \ja{きたら、かぞくの} \ja{あんぴを} \ja{かくにんします。}} — \emph{taifū ga kitara, kazoku no anpi o kakunin shimasu} — \textquotedblleft{}If a typhoon comes, I check that my family are safe.\textquotedblright{}
\end{quote}

\begin{cousinweb}
\textbf{\ja{あんぴ}} (\emph{anpi}) is Sino-Japanese (kango): two Chinese characters read with their on'yomi, readings borrowed from Chinese. They are \emph{an}, \textquotedblleft{}peaceful\textquotedblright{} (Middle Chinese \emph{'an}) and \emph{hi}, \textquotedblleft{}negate\textquotedblright{} (KANJIDIC2; Middle Chinese in Baxter's transcription). After \emph{an} the second is said \emph{pi}.
\end{cousinweb}

\begin{grammarlens}[title={What you've built}]
One more word for saying what happens if.
\end{grammarlens}

\subsection*{Your turn}
\begin{itemize}
\raggedright
  \item \emph{Say it:} \emph{anpi}
  \item \emph{Say it:} \emph{anpi}, once more
  \item \emph{Say it:} \emph{taifū ga kitara, kazoku no anpi o kakunin shimasu} — If a typhoon comes, I check that my family are safe.
  \item \emph{Recall:} say the Japanese for insurance, then the Japanese for to confirm, then say \emph{anpi} again
\end{itemize}

\begin{tcolorbox}[breakable,colback=teal!4,colframe=teal!35!black,title={Before you move on}]
What does \ja{あんぴ} mean? (\textquotedblleft{}Whether someone is safe\textquotedblright{}.) Say it once more.
\end{tcolorbox}
\section[\texorpdfstring{\ja{ねんのため} (nen no tame)}{ねんのため (nen no tame)}]{\ja{ねんのため} (nen no tame) — to be on the safe side; to make sure}

\label{lesson:JA-C351-nen-no-tame}



\begin{quote}
Before the new one: say the Japanese for to confirm, then the Japanese for whether someone is safe.
\end{quote}

\subsection*{You'll want to know: \ja{ねんのため}}
\textbf{\ja{ねんのため}} — \emph{nen no tame} — \textquotedblleft{}to be on the safe side; to make sure\textquotedblright{}.

\begin{quote}
\textbf{\ja{ねんのため、ついたら} \ja{でんわして} \ja{ください。}} — \emph{nen no tame, tsuitara denwa shite kudasai} — \textquotedblleft{}To be on the safe side, please phone when you arrive.\textquotedblright{}
\end{quote}

\begin{cousinweb}
\textbf{\ja{ねんのため}} (\emph{nen no tame}) is built from the Sino-Japanese \emph{nen}, \textquotedblleft{}sense\textquotedblright{} or attention, an on'yomi borrowed from Chinese (Middle Chinese \emph{nemH}), and the native \emph{no tame}, \textquotedblleft{}for the sake of\textquotedblright{}: \textquotedblleft{}for the sake of care\textquotedblright{} (KANJIDIC2; JMdict; Middle Chinese in Baxter's transcription).
\end{cousinweb}

\begin{grammarlens}[title={What you've built}]
One more word for saying what happens if.
\end{grammarlens}

\subsection*{Your turn}
\begin{itemize}
\raggedright
  \item \emph{Say it:} \emph{nen no tame}
  \item \emph{Say it:} \emph{nen no tame}, once more
  \item \emph{Say it:} \emph{nen no tame, tsuitara denwa shite kudasai} — To be on the safe side, please phone when you arrive.
  \item \emph{Recall:} say the Japanese for to confirm, then the Japanese for whether someone is safe, then say \emph{nen no tame} again
\end{itemize}

\begin{tcolorbox}[breakable,colback=teal!4,colframe=teal!35!black,title={Before you move on}]
What does \ja{ねんのため} mean? (\textquotedblleft{}To be on the safe side; to make sure\textquotedblright{}.) Say it once more.
\end{tcolorbox}
\section[(dialogue)]{Looking back: from Visiting Someone to In Case Something Happens}

\label{lesson:JA-R351-review-347-351}



\begin{quote}
Say the Japanese for to pay a visit and the matter in hand.
\end{quote}

\subsection*{Your turn}
Say these aloud:

\begin{itemize}
\raggedright
  \item to pay a visit, the matter in hand, in person, while you are at it, a caller
  \item to be absent, slow, at the last minute, the hour, consecutive holidays
  \item an expectation, the situation, an opportunity, to miss a chance, at great pains
  \item to take precautions, carelessness, to prepare for, a countermeasure, disaster preparedness
  \item in the unlikely event, insurance, to confirm, whether someone is safe, to be on the safe side
\end{itemize}

\begin{tcolorbox}[breakable,colback=teal!4,colframe=teal!35!black,title={Before you move on}]
To pay a visit? (\textbf{\ja{ほうもんする}}, \emph{hōmon suru}.) The matter in hand? (\textbf{\ja{ようけん}}, \emph{yōken}.) In person? (\textbf{\ja{ちょくせつ}}, \emph{chokusetsu}.) While you are at it? (\textbf{\ja{ついでに}}, \emph{tsuideni}.) A caller? (\textbf{\ja{らいきゃく}}, \emph{raikyaku}.) To be absent? (\textbf{\ja{けっせきする}}, \emph{kesseki suru}.) Slow? (\textbf{\ja{おそい}}, \emph{osoi}.) At the last minute? (\textbf{\ja{ぎりぎり}}, \emph{girigiri}.) The hour? (\textbf{\ja{じこく}}, \emph{jikoku}.) Consecutive holidays? (\textbf{\ja{れんきゅう}}, \emph{renkyū}.) An expectation? (\textbf{\ja{よそう}}, \emph{yosō}.) The situation? (\textbf{\ja{じょうきょう}}, \emph{jōkyō}.) An opportunity? (\textbf{\ja{きかい}}, \emph{kikai}.) To miss a chance? (\textbf{\ja{のがす}}, \emph{nogasu}.) At great pains? (\textbf{\ja{せっかく}}, \emph{sekkaku}.) To take precautions? (\textbf{\ja{ようじんする}}, \emph{yōjin suru}.) Carelessness? (\textbf{\ja{ゆだん}}, \emph{yudan}.) To prepare for? (\textbf{\ja{そなえる}}, \emph{sonaeru}.) A countermeasure? (\textbf{\ja{たいさく}}, \emph{taisaku}.) Disaster preparedness? (\textbf{\ja{ぼうさい}}, \emph{bōsai}.) In the unlikely event? (\textbf{\ja{まんいち}}, \emph{man'ichi}.) Insurance? (\textbf{\ja{ほけん}}, \emph{hoken}.) To confirm? (\textbf{\ja{かくにんする}}, \emph{kakunin suru}.) Whether someone is safe? (\textbf{\ja{あんぴ}}, \emph{anpi}.) To be on the safe side? (\textbf{\ja{ねんのため}}, \emph{nen no tame}.)
\end{tcolorbox}
